\section{All tensor-map derivatives and the complementary Capelli operators}
\label{sec:bessel-shift-suppliers-v51}

This chapter supplies the whole original \texttt{TMapStatement} and
\texttt{TMapDStatement}, the two complementary Capelli identities, and
\texttt{EigStatement}.  In particular, the estimates apply to every
tuple from the generated module and to every derivative word, with no
loss in the original exponent.  The final claim isolates the remaining
period-evaluation and shifted-quotient inputs to the two unit shifts.
Their existence is not inferred from the algebra below.

We use the reviewed DC, BNV and BSP notation.  Thus
$G=U(J_\epsilon)$, $K=U(p)\times U(q)$, $\rho(g)=\|g\|_{\rm F}$,
$q_0=\min(p,q)$, and
\[
 a=\frac{h}{8\sqrt{2q_0}},\qquad \beta=a-1,\qquad
 \Lambda=\lambda+\rho_n-\rho_c.
 \tag{BST.1}
\]
Statements about arbitrary dimension or arbitrary algebra below do not
impose the Bessel signature.  In the tensor-map application
$n=2r+k+1$, $p=k+r$, $q=r+1$, and the original \texttt{StepGap}
is assumed.  BSP1 then gives $q_0\ge1$, $h>h_{\rm need}$, and
$\beta>4\max(2nq_0,2rn^2)$.

Our inner product is conjugate-linear in its first slot.  Write
$R_XF(g)=\frac{d}{dt}|_{t=0}F(g\exp(tX))$ for a real Lie direction.
The source's $E_{ab}$ on functions is its complex-linear extension to
the matrix unit.  For a word $\mathbf w=(w_1,\ldots,w_\ell)$,
$E_{\mathbf w}=E_{w_1}\cdots E_{w_\ell}$, with the rightmost
operator applied first.  Smoothness on $G$ always means the original
local extension to a smooth real function on ambient matrices.
No claim requires the originally given off-group values to be smooth.

\subsection{Compact smoothing without loss of decay}

% BEGIN BST1
\paragraph{BST1. A compact smooth convolution fixing the entire lowest type.}
\label{claim:BST1}
For arbitrary $n,\epsilon$ and a kernel datum
$(\lambda,h,T,\mu)$ at the original
Harish--Chandra, gap, $q_0\ge1$, $h>h_{\rm need}$ and
$\mu=\lambda+\rho$ hypotheses, there exists one smooth compactly
supported function $b$ on ambient complex matrices such that
\[
 \Psi_{T,\mu}(g)
    =\int_G b(y)\Psi_{T,\mu}(gy)\,dy\qquad(g\in G).
 \tag{BST.2}
\]
The Haar normalization is arbitrary.  The same $b$ works for every
pair of vectors in $T$.

\begin{proof}
Take the normalized Hilbert model of DC1--DC11, writing $j:T\to
\mathcal H$ for its isometric lowest-type embedding.  Its coefficients
are exactly $\Psi$ and the full $T$-isotypic subspace is $j(T)$.
The latter assertion is the multiplicity and closure conclusion of
DC9--DC10, not an additional admissibility assumption.

Choose a nonnegative smooth approximate identity $b_0$ on $G$, of
compact support and integral one.  Average it under conjugation by
$K$:
\[
 b_1(y)=\int_K b_0(k^{-1}yk)\,dk ,
 \tag{BST.3}
\]
where compact Haar has mass one.  This is smooth, and its support is
contained in the compact set $K(\operatorname{supp}b_0)K^{-1}$.
Conjugation by compact elements preserves group Haar measure: its
modulus is a continuous homomorphism from a compact group to
$\mathbb R_{>0}$ and hence equals one.
The bounded operator $A=\int_G b_1(y)\sigma(y)\,dy$ commutes with
$K$.  It consequently carries $j(T)$ into that same full isotypic
subspace, and Schur's lemma makes its restriction a scalar $c$.

By strong continuity, the operators obtained from a sufficiently small
$b_0$ are arbitrarily close to the identity on this finite-dimensional
subspace, uniformly on its unit sphere.  Averaging the conjugates
preserves that estimate, because $K$ preserves $j(T)$ isometrically.
Choose the support small enough that $|c-1|<1/2$.  Then $c\ne0$,
and $b=b_1/c$ satisfies
$\int_G b(y)\sigma(y)j(v)\,dy=j(v)$ for every $v\in T$.
Pairing after $\sigma(g)$ with $j(v')$ proves (BST.2), first as
scalar coefficients and hence as an equality of finite-dimensional
endomorphisms.  Every integral here is absolutely convergent because
its integrand has compact support and the representation is unitary.

Finally, a compactly supported smooth function on the closed embedded
submanifold $G$ has a compactly supported smooth ambient extension.
Indeed cover its compact support by finitely many submanifold charts,
use a smooth partition of unity, extend each chart function constantly
in the normal coordinates, and multiply by a compact bump function
equal to one near its support.  Sum these finitely many extensions.
Replacing $b$ by this extension changes no integral in (BST.2).
A change of Haar normalization is absorbed by the inverse scaling
of $b$.
\end{proof}
% END BST1

% BEGIN BST2
\paragraph{BST2. The whole original \texttt{DecDerivStatement}.}
\label{claim:BST2}
At every datum of BST1, for every matrix-unit word $\mathbf w$
and every $v_1,v_2\in T$ there exists a constant $C_{\mathbf w,v_1,v_2}$
such that
\[
 \left|E_{\mathbf w}
       \langle v_2,\Psi_{T,\mu}(g)v_1\rangle\right|
       \le C_{\mathbf w,v_1,v_2}\rho(g)^{-a}
       \qquad(g\in G).
 \tag{BST.4}
\]
The constant may depend on the word and vectors; the exponent does not.

\begin{proof}
We first recall explicitly how DC6 gives the required undifferentiated
Frobenius bound at this arbitrary signature.  Write $g=ke^X$,
$s=\|X\|_{\rm F}$ and $R=\sqrt{2q_0}\ge1$.  DC6 gives
\[
 \|\Psi(g)\|\le
  \exp\left[-\frac h8\min(s^2,s/R)\right],
 \qquad \rho(g)\le\sqrt n\,e^s.
\]
The elementary inequality
$\min(s^2,s/R)\ge(s-1)/R$ holds: if $sR\ge1$ the minimum
is $s/R$, and otherwise $s<1$ and the right side is nonpositive.
Thus
\[
           \|\Psi(g)\|\le e^a n^{a/2}\rho(g)^{-a}.
 \tag{BST.5}
\]
Here $a>0$, since the displayed large-gap hypothesis implies $h>0$.

Let $\phi(g)=\langle v_2,\Psi(g)v_1\rangle$.
By BST1, $\phi(g)=\int_G b(y)\phi(gy)\,dy$.
For a real Lie direction $X$, left Haar invariance gives
\[
 \phi(g\exp(tX))
    =\int_G b(\exp(-tX)y)\phi(gy)\,dy.
 \tag{BST.6}
\]
For $t$ in a fixed small interval the supports on the right lie in
one compact set.  All derivatives of the smooth factor $b$ are
bounded there.  For $g$ in a fixed compact neighborhood, continuity
of $\phi$ bounds the other factor as well.  Dominated differentiation
therefore applies to every order, locally uniformly in $g$.
In particular, set
$D_Xb(y)=\frac{d}{dt}|_{0}b(\exp(tX)y)$; the kernel for $R_X\phi$
is $-D_Xb$.  Extend this operation complex linearly to each
matrix unit and iterate in the specified word order.  We obtain a
smooth function $b_{\mathbf w}$ with support contained in
$\operatorname{supp}(b|_G)$ and
\[
 E_{\mathbf w}\phi(g)=
           \int_G b_{\mathbf w}(y)\phi(gy)\,dy .
 \tag{BST.7}
\]
The same construction may be performed on the ambient extension of
$b$; its directional differential operators have smooth coefficients
and preserve compact support.  Thus it also has the source's precise
ambient meaning.

On that fixed compact support, let $B$ bound $\rho(y^{-1})$.
Frobenius submultiplicativity gives
$\rho(g)\le\rho(gy)\rho(y^{-1})$, hence
$\rho(gy)^{-a}\le B^a\rho(g)^{-a}$.
Apply (BST.5) inside (BST.7) and integrate
$|b_{\mathbf w}|$:
\[
 |E_{\mathbf w}\phi(g)|
 \le e^a n^{a/2}\|v_2\|\|v_1\|\,
       B^a\|b_{\mathbf w}\|_{L^1}\rho(g)^{-a}.
\]
This is (BST.4).  It proves the complete source statement for all
words, including the empty word.  No estimate for a differentiated
kernel was assumed.
\end{proof}
% END BST2

% BEGIN BST3
\paragraph{BST3. Every generated-module coefficient, with all words and translates.}
\label{claim:BST3}
Take any of the original complete unitary derived data
$(\mathcal H,\sigma,D,E,V,T,\mu,\iota)$ in
\texttt{TMapStatement}.  In particular its lowest-type coefficients
equal the specified $\Psi$ kernel, its derived domain is invariant,
and all the original gap hypotheses hold.
Put $S=U(\mathfrak{gl}_n)V\subset D$.
For every $x_0\in G$, every $w\in V$ and every $m\in S$, the
function
\[
 F_m(g)=\langle w,\sigma(x_0^{-1}g)m\rangle
 \tag{BST.8}
\]
is smooth in the original local sense.  For every word $\mathbf w$,
\[
 E_{\mathbf w}F_m=F_{E_{\mathbf w}m},\qquad
 |E_{\mathbf w}F_m(g)|\le C_{m,\mathbf w,x_0,w}\rho(g)^{-a}.
 \tag{BST.9}
\]
No irreducibility of the ambient representation is required.

\begin{proof}
The algebraic subspace spanned by all finite words $E_{\boldsymbol\alpha}v$,
$v\in V$, contains $V$ and is stable under every $E_{ab}$.
Conversely every such stable subspace contains every word and their
finite linear combinations.  It is therefore exactly the source's
generated module $S$.

For $m=v\in V$, the original coefficient identity makes (BST.8)
equal on $G$ to
$\langle\iota w,\Psi(x_0^{-1}g)\iota v\rangle$.
Its smoothness follows from DC2.  Right differentiation commutes with
the fixed left translate $x_0^{-1}$.
The defining derived-representation identity gives
$E_{ab}F_m=F_{E_{ab}m}$, first for real directions by differentiation
of $\sigma(x_0^{-1}g)\sigma(\exp(tX))m$, and then by complexification.
Iterating gives the first formula of (BST.9).
For $m=E_{\boldsymbol\alpha}v$, the resulting coefficient is a
derivative of the preceding smooth kernel coefficient; its further
$\mathbf w$ derivative is the kernel derivative for the concatenated
word $\mathbf w\boldsymbol\alpha$.
BST2 bounds it by a constant times $\rho(x_0^{-1}g)^{-a}$.
Since $\rho(g)\le\rho(x_0)\rho(x_0^{-1}g)$, this is at most a
constant times $\rho(g)^{-a}$.  The constant $\rho(x_0)^a$ is
independent of $g$.
Taking finite linear combinations proves both assertions for every
$m\in S$.  All the local smooth extensions used here may be chosen
after the fixed translate; on the group they are the actual functions
(BST.8), so equality of tangential derivatives is unaffected by
off-group values.
\end{proof}
% END BST3

\subsection{The paired tensor maps at their original scope}

% BEGIN BST4
\paragraph{BST4. Exact standard and dual tensor intertwining.}
\label{claim:BST4}
For every $x_1\in\mathbb C^n$, every $x_0\in G$, $w\in V$, and
every tuple $f=(f_j)_{j=0}^{n-1}\in S^n$, define the original maps
\[
 \begin{aligned}
 Q_f(g)&=\sum_j q_j(g)F_{f_j}(g),
       &q_j(g)&=x_1^*J_\epsilon ge_j,\\
 Q_f^\vee(g)&=\sum_j r_j(g)F_{f_j}(g),
       &r_j(g)&=e_j^{\,t}g^{-1}x_1.
 \end{aligned}
 \tag{BST.10}
\]
In particular $x_1$ is arbitrary; it need not come from a Bessel frame.
The two functions are smooth, and the exact tensor actions on tuples are
\[
 \begin{split}
 (L_{ab}^{+}f)_c&=E_{ab}f_c+\delta_{ca}f_b,\\
 (L_{ab}^{-}f)_c&=E_{ab}f_c-\delta_{cb}f_a.
 \end{split}
 \tag{BST.11}
\]
They satisfy on $G$
\[
                  E_{ab}Q_f=Q_{L_{ab}^{+}f},
 \qquad E_{ab}Q_f^\vee=Q_{L_{ab}^{-}f}^\vee .
 \tag{BST.12}
\]
The actions in (BST.11) are the literal \texttt{LemK.tensV}
and \texttt{LemK.tensVd}.

\begin{proof}
The first multiplier is linear in the matrix entries; the second is
holomorphic on invertible matrices.  Both have the required local
real smoothness.  Differentiating $g\exp(tX)$ and its inverse gives
\[
 E_{ab}q_j=\delta_{jb}q_a,\qquad
 E_{ab}r_j=-\delta_{ja}r_b.
 \tag{BST.13}
\]
These are first obtained along real Lie directions and then extended
complex linearly, so their conventions agree with the original
$E_{ab}$.  By BST3 the coefficient factors are smooth and their
derivatives act on $f_j$.
The ordinary product rule now gives
\[
 E_{ab}Q_f=\sum_jq_jF_{E_{ab}f_j}+q_aF_{f_b},\qquad
 E_{ab}Q_f^\vee=\sum_jr_jF_{E_{ab}f_j}-r_bF_{f_a}.
 \tag{BST.14}
\]
Linearity of $F_m$ in $m$ turns these formulas exactly into
(BST.11)--(BST.12).  All tuples on the right still lie in $S^n$,
since $S$ is a submodule stable under every $E_{ab}$.
\end{proof}
% END BST4

% BEGIN BST5
\paragraph{BST5. The whole original \texttt{TMapStatement} and \texttt{TMapDStatement}.}
\label{claim:BST5}
Under every original datum and hypothesis of those two statements,
for every $x_1,x_0,w,f$ as in BST4 and every derivative word
$\mathbf w$ there is a real constant $C_{\mathbf w}$ such that
\[
 \begin{aligned}
 E_{\mathbf w}Q_f&=Q_{L_{\mathbf w}^{+}f},
 &|E_{\mathbf w}Q_f(g)|&\le C_{\mathbf w}\rho(g)^{1-a},\\
 E_{\mathbf w}Q_f^\vee&=Q_{L_{\mathbf w}^{-}f}^\vee,
 &|E_{\mathbf w}Q_f^\vee(g)|&\le C_{\mathbf w}\rho(g)^{1-a}.
 \end{aligned}
 \tag{BST.15}
\]
The two constants may be enlarged to one as displayed.
Both functions and every derivative word lie in the actual
$L^2(G,\mu_G)$ of the source.  Thus both original
\texttt{IsSmoothL2} conclusions, both exact first-derivative
intertwinings, and all original word estimates hold.

\begin{proof}
Iterate BST4.  Define
$L_\varnothing^\pm f=f$ and
$L_{(a,b)::\mathbf w}^\pm f=L_{ab}^\pm L_{\mathbf w}^\pm f$.
Induction using (BST.11) keeps every component of every such tuple
in $S$ and proves the equalities in (BST.15), including their order.

There is one constant $B_{x_1}$ such that, for all $j$ and $g\in G$,
\[
                    |q_j(g)|,\ |r_j(g)|\le B_{x_1}\rho(g).
 \tag{BST.16}
\]
For $q_j$ this follows directly from Cauchy--Schwarz and
$\|J_\epsilon\|_{\rm op}=1$.  For $r_j$ use
$g^{-1}=J_\epsilon g^*J_\epsilon$, so
$\rho(g^{-1})=\rho(g)$.
Fix a word and apply BST3 to each component of
$L_{\mathbf w}^\pm f$.  Sum the finitely many coefficient bounds
and multiply by (BST.16).  This proves the exact exponent $1-a$,
with no extra factor depending on word length in the exponent.

Finally $\beta=a-1>0$ and $2\beta>2nq_0$ by BSP1.
The square of either bound in (BST.15) is a constant times
$\rho^{-2\beta}$, which is integrable by BSP6 for the actual Haar
normalization.  Smoothness gives measurability, so every derivative
word is genuinely in $L^2$, rather than merely having a formal
\texttt{toL2} value.  Smoothness itself was proved in BST4 and is
preserved by every tangential derivative.  These are precisely all
fields of the source's \texttt{IsSmoothL2}.
The argument retains arbitrary generated-module tuples, arbitrary
free $x_0$, all $x_1$, every lowest-type $w$, and all original
signature and gap assumptions.
\end{proof}
% END BST5

\subsection{Complementary minors over an arbitrary complex algebra}

For this subsection let $\mathcal U$ be any associative unital complex
algebra and let $E_{ij}\in\mathcal U$ obey
\[
 [E_{ij},E_{kl}]=\delta_{jk}E_{il}-\delta_{li}E_{kj}.
 \tag{BST.17}
\]
Put $A_{ij}(t)=E_{ij}+t\delta_{ij}$.  For lists
$R,C:\{0,\ldots,p-1\}\to\{0,\ldots,n-1\}$, define
\[
 \begin{split}
 C_{R,C}(x)
  &=\sum_{\sigma\in S_p}\operatorname{sgn}(\sigma)
       \prod_{j=0}^{p-1}A_{R_{\sigma(j)},C_j}(x-j),\\
 D_{R,C}(x)
  &=\sum_{\sigma\in S_p}\operatorname{sgn}(\sigma)
       \prod_{i=0}^{p-1}A_{R_i,C_{\sigma(i)}}(x-p+1+i).
 \end{split}
 \tag{BST.18}
\]
Products have the displayed left-to-right order.  These are exactly
the source's column and row minors; empty products equal one.
A hat over an index deletes that position and retains the increasing
order of all remaining positions.

% BEGIN BST6
\paragraph{BST6. Alternation and the exact minor commutator.}
\label{claim:BST6}
Both minors in (BST.18) are alternating in both index lists, including
vanishing if either list repeats an index.  Moreover
\[
 [E_{ab},C_{R,C}(x)]
 =\sum_{i:R_i=b}C_{R[i\leftarrow a],C}(x)
  -\sum_{j:C_j=a}C_{R,C[j\leftarrow b]}(x).
 \tag{BST.19}
\]
All assertions hold in every algebra satisfying (BST.17).

\begin{proof}
Row alternation of $C$ and column alternation of $D$ are immediate
by reindexing their defining permutation sums.  To prove the other
two alternations without commuting noncommuting entries, adjoin
exterior variables $e_0,\ldots,e_{n-1}$ commuting with $\mathcal U$.
Set
$\eta_j=\sum_i e_iE_{ij}$ and $\eta_j(t)=\eta_j+te_j$.
Directly from (BST.17),
\[
 \eta_i\eta_j+\eta_j\eta_i=-e_i\eta_j-e_j\eta_i,
 \qquad
 \eta_i(t)\eta_j(t-1)=-\eta_j(t)\eta_i(t-1).
 \tag{BST.20}
\]
For example the first identity follows by writing its left side as
$\sum_{a,b}e_ae_b[E_{ai},E_{bj}]$; the two Kronecker sums give
the two displayed terms.  Expanding the scalar shifts proves the
second identity.  The coefficient of
$e_{R_0}\cdots e_{R_{p-1}}$ in
$\eta_{C_0}(x)\cdots\eta_{C_{p-1}}(x-p+1)$ is
$C_{R,C}(x)$ when the rows are distinct, with the indicated ordering.
Adjacent column interchange negates the whole product by (BST.20).
Together with the already proved row alternation, this gives column
alternation for every row list as well.

For the row version set $\xi_i=\sum_j e_jE_{ij}$.
The same calculation gives
$\xi_i\xi_j+\xi_j\xi_i=e_i\xi_j+e_j\xi_i$, and therefore
\[
       \xi_i(t)\xi_j(t+1)=-\xi_j(t)\xi_i(t+1).
 \tag{BST.21}
\]
Coefficients in
$\xi_{R_0}(x-p+1)\cdots\xi_{R_{p-1}}(x)$ prove row alternation of
$D$.  A repeated list entry gives its own negative under a
transposition; characteristic zero makes the minor zero.

Finally
$[E_{ab},A_{cd}(t)]=
\delta_{bc}A_{ad}(t)-\delta_{da}A_{cb}(t)$.
The scalar correction terms on the right cancel each other.
Apply the Leibniz rule for a commutator to each ordered product
in $C$.  In the first set of terms, replace its unique row position
by $a$ and sum over that position; in the second set replace the
specified column position by $b$.  Summing over permutations is
exactly (BST.19).  No reordering of the product factors occurs.
\end{proof}
% END BST6

% BEGIN BST7
\paragraph{BST7. Equality of the shifted row and column minors.}
\label{claim:BST7}
For every $p,R,C,x$,
\[
                         C_{R,C}(x)=D_{R,C}(x).
 \tag{BST.22}
\]
In particular the first-row expansion of the column minor is
\[
 C_{R,C}(x)=
 \sum_{j=0}^{p}(-1)^j A_{R_0,C_j}(x-p)
           C_{R\widehat{0},C\widehat{j}}(x)
 \tag{BST.23}
\]
when $R,C$ have length $p+1$.

\begin{proof}
Induct on the common list length $N$.  The empty case is one.
Assume the assertion for length $N-1$, and abbreviate
$C_{ij}=C_{R\widehat{i},C\widehat{j}}(x-1)$.
Ordinary first-column expansion and column alternation from BST6
allow any chosen column to be moved to the first position; hence
\[
 N C_{R,C}(x)=
       \sum_{i,j}(-1)^{i+j}A_{R_i,C_j}(x)C_{ij}.
 \tag{BST.24}
\]
Similarly, use the last-row expansion of the row minor, moving each
chosen row to the last position.  Its entry at that position has
parameter $x$.  The induction hypothesis on the remaining minor
gives
\[
 N D_{R,C}(x)=
       \sum_{i,j}(-1)^{i+j}C_{ij}A_{R_i,C_j}(x).
 \tag{BST.25}
\]
These expansions are valid also for repeated indices, by alternation.

We check the entire commutator difference of the two double sums.
Define
\[
 T_{sj}=\delta_{R_s,C_j}(-1)^{s+j+1}
                  C_{R\widehat{s},C\widehat{j}}(x-1).
 \tag{BST.26}
\]
Formula (BST.19) gives exactly
\[
 (-1)^{i+j}[E_{R_i,C_j},C_{ij}]
               =\sum_{s\ne i}T_{sj}-\sum_{t\ne j}T_{it}.
 \tag{BST.27}
\]
To verify its signs, a term in the first commutator sum requires
$R_s=C_j$, $s\ne i$.  Replacing row $R_s$ in the list with $i$
deleted by $R_i$, and putting the result into the order with $s$
deleted, has sign $(-1)^{i+s+1}$.  Multiplication by
$(-1)^{i+j}$ leaves $(-1)^{s+j+1}$, the sign of $T_{sj}$.
The column replacements similarly leave $(-1)^{i+t+1}$ and occur
with the minus sign in (BST.19).  This proves (BST.27).
The scalar part of $A_{R_i,C_j}(x)$ commutes and contributes nothing.

On summing (BST.27) over $i,j$, every $T_{sj}$ occurs $N-1$
times in the first sum and $N-1$ times in the second.  The
difference is zero.  Equations (BST.24)--(BST.25) therefore give
$N(C_{R,C}-D_{R,C})=0$; division by the nonzero complex integer
$N$ completes the induction.  Finally apply ordinary first-row
expansion to $D$ and (BST.22) to its remaining minors.  Its first
parameter is $x-p$, which proves the exact shift in (BST.23).
\end{proof}
% END BST7

% BEGIN BST8
\paragraph{BST8. The full original complementary Capelli expansions.}
\label{claim:BST8}
Let $n=m+1$, $\mathcal C(x)=C_{\mathrm{id},\mathrm{id}}(x)$,
and define
\[
 U_{dc}(x)=(-1)^{d+c}C_{\widehat d,\widehat c}(x-1),
 \qquad
 V_{dc}(x)=(-1)^{d+c}C_{\widehat c,\widehat d}(x).
 \tag{BST.28}
\]
Then for every complex $x$,
\[
 \begin{split}
 \sum_d(E_{db}+x\delta_{db})U_{dc}(x)
                         &=\delta_{bc}\mathcal C(x),\\
 \sum_d(E_{ad}+(x-m)\delta_{ad})V_{dc}(x)
                         &=\delta_{ac}\mathcal C(x).
 \end{split}
 \tag{BST.29}
\]
This is the complete original \texttt{CapelliExpandStatement}.

\begin{proof}
For the first identity take the column minor with full row list and
column list $(b,0,\ldots,\widehat c,\ldots,m)$.
Expand its first column.  The terms are
$(-1)^d(E_{db}+x\delta_{db})
C_{\widehat d,\widehat c}(x-1)$.
Multiply their sum by $(-1)^c$.  If $b=c$, the column list is
the identity list with $c$ moved to the front, so its sign is
$(-1)^c$.  If $b\ne c$, the list repeats $b$ and the minor
vanishes by BST6.  This proves the first identity, including the
order of every factor.

For the second take row list $(a,0,\ldots,\widehat c,\ldots,m)$
and full column list.  Use (BST.23), whose first-row entries
have parameter $x-m$, and again multiply by $(-1)^c$.
Row alternation gives the displayed right side.  Empty complementary
minors when $m=0$ are one, so no positive-rank exception was used.
\end{proof}
% END BST8

% BEGIN BST9
\paragraph{BST9. The whole original \texttt{EigStatement}.}
\label{claim:BST9}
Let $M$ be any complex module, $E_{ij}\in\operatorname{End}_{\mathbb C}M$
satisfy (BST.17), and $S\subset M$ be stable under every $E_{ij}$.
Let $m\ge0$ and suppose, for one real $\omega$, that
$\mathcal C(\omega)s=0$ for every $s\in S$.
For arbitrary $y,\alpha\in\mathbb C^{m+1}$ and $v\in S$, set
\[
 (n^+)_d=\sum_c y_cU_{dc}(\omega)v,\qquad
 (n^-)_d=\sum_c \alpha_cV_{dc}(\omega)v.
 \tag{BST.30}
\]
These are exactly \texttt{nnPlus} and \texttt{nnMinus}.
Writing
\[
 (\Omega^+f)_b=\sum_d E_{db}f_d,\qquad
 (\Omega^-f)_a=-\sum_d E_{ad}f_d,
 \tag{BST.31}
\]
one has
\[
 n^+\in S^{m+1}\cap\ker(\Omega^++\omega),\qquad
 n^-\in S^{m+1}\cap\ker(\Omega^--(\omega-m)).
 \tag{BST.32}
\]
Thus the eigenvalues are the original $-\omega$ and $\omega-m$.

\begin{proof}
Every entry in a complementary minor is a matrix-unit operator plus
a scalar.  Stability of $S$, followed by finite products and sums,
puts every component of (BST.30) in $S$.
Apply the first identity of (BST.29) to $v$, multiply by $y_c$,
and sum over $c$.  Its right side is zero.  The left side is
$(\Omega^+n^+)_b+\omega(n^+)_b$, giving the first kernel condition.
The second identity, multiplied by $\alpha_c$ and summed, gives
$\sum_dE_{ad}(n^-)_d+(\omega-m)(n^-)_a=0$.
The minus sign in the definition of $\Omega^-$ gives the second
condition exactly.  These proofs work over an arbitrary module;
they assume no Hilbert structure, nonzero vector, or irreducibility.
The complex version of the two eigenvalue identities also holds;
real $\omega$ is retained to match the source's real eigenvalue
parameter in \texttt{Xs}.
\end{proof}
% END BST9

% BEGIN BST10
\paragraph{BST10. The full tensor Casimir identities and both exact shifts.}
\label{claim:BST10}
For every $n,M,E$ as above, the two actions $L^\pm$ of (BST.11)
are representations of $\mathfrak{gl}_n$ and obey
\[
 C_{L^\pm}f=(C_Ef_j)_j+n f+2\Omega^\pm f,
 \qquad C_E=\sum_{a,b}E_{ab}E_{ba}.
 \tag{BST.33}
\]
This is the complete original \texttt{CapelliTensorStatement}.
In the original kernel application, put
$\theta_s=\lambda_s-(n-1)/2$ and $\omega=-\theta_s$.
BSP3 and BST9 then give the prescribed standard and dual
eigenspace vectors.  On those eigenspaces their tensor Casimirs are
respectively
\[
                 c_{\lambda+e_s},\qquad c_{\lambda-e_s}.
 \tag{BST.34}
\]
The vectors are not asserted nonzero.

\begin{proof}
Write $L^\pm_{ab}$ as the componentwise operator $E_{ab}$ plus
the standard, respectively contragredient, matrix-unit action on
the tuple index.  The two summands commute.  Each satisfies the
matrix-unit bracket relations, so their sum does also.
Expand $\sum_{a,b}L^\pm_{ab}L^\pm_{ba}$.
The componentwise terms give $(C_Ef_j)_j$.
On either the standard or dual index representation,
$\sum_{a,b}e_{ab}e_{ba}=nI$; in the dual action the two minus
signs cancel.  The two cross terms agree after exchanging $a,b$.
For the standard action each has component
$\sum_dE_{db}f_d$; for the dual each has component
$-\sum_dE_{ad}f_d$.  They are precisely the two terms
$2\Omega^\pm$ in (BST.33).

In the application $C_E=c_\lambda$ on $S$.
Indeed it has that scalar on $V$ by the original datum and commutes
with every $E_{ab}$ by the matrix-unit relations, so the identity
extends to the generated module.
The eigenvalues from BST9 are
$\theta_s$ and $-\theta_s-(n-1)$, respectively.
Since $c_\lambda=\|\lambda\|^2-\|\rho\|^2$,
\[
 c_\lambda+n+2\theta_s=c_\lambda+2\lambda_s+1
        =c_{\lambda+e_s},\qquad
 c_\lambda+n+2(-\theta_s-(n-1))
        =c_\lambda-2\lambda_s+1=c_{\lambda-e_s}.
 \tag{BST.35}
\]
This proves the literal shifts, including the contragredient sign.
It does not identify a lowest compact type in an image or quotient.
\end{proof}
% END BST10

\subsection{The exact residual unit-shift problem}

% BEGIN BST11
\paragraph{BST11. Reduction of the whole original \texttt{StepStatement} to four precise suppliers.}
\label{claim:BST11}
Retain the full original \texttt{EvalPStatement},
\texttt{EvalMStatement}, \texttt{QuotPStatement} and
\texttt{QuotMStatement}, at the exact data and quantifiers in
\texttt{LaplaceEstimates/PsiDatumStatement}.  Then the reviewed
DC/BNV/BSP results and BST1--BST10 prove the whole original
\texttt{StepStatement}.
Thus these four suppliers, rather than a derivative-decay or
complementary-eigenspace input, are the remaining assumptions in
this analytic reduction.

\begin{proof}
Fix arbitrary original $k,r,\epsilon,u,h$ with the signature and
\texttt{StepGap} hypotheses.  The original subgroup-integrability
antecedent is supplied by BNV1.
Take a witness to $\texttt{BesselNZ}(\epsilon,\lambda,u)$ and
one of the two allowed endpoints $\lambda\pm e_s$.
The definition of an allowed step retains its Harish--Chandra and
noncompact-gap hypotheses, as well as the appropriate exclusion
$u_j\ne\lambda_s\pm1/2$ for every $j$.

By BSP1 and BSP2 choose an exact kernel datum for the old
$T,\mu$.  If the old left translate is $x_0\in G_H$, put
$x'_0=Px_0P^{-1}\in G$.  Under the isomorphism $\iota$,
the two old kernel vectors become $v,w\in V$.
On the actual Haar support the coefficient identity is
\[
 \langle w,\sigma((x'_0)^{-1}P h'n'xP^{-1})v\rangle
 =
 \langle v_2,\Psi(P(x_0^{-1}h'n'x)P^{-1})v_1\rangle.
 \tag{BST.36}
\]
Hence its original \texttt{per} integral is the old nonzero
\texttt{aaFun} integral.  The Cayley matrix, character, compact
representation, Haar measures, and the vectors and translates
are exactly the old ones.

For the positive shift, the retained \texttt{EvalP} supplies $y$
such that the period of the exact tensor map of
$\texttt{nnPlus}(y,-\theta_s,v)$ is nonzero.
BSP3 kills $\mathcal C(-\theta_s)$ on all $S$; BST9 places this
tuple in the original standard $\theta_s$ eigenspace.
Apply the retained \texttt{QuotP} to this tuple and nonzero period.
It supplies the actual spaces $W_f,W'$ with their exact image
relation, shifted lowest type and Casimir, actual compact stability,
and period injectivity.
Every $F\in W_f$ is a tensor map of a tuple in this eigenspace.
BST5 gives its continuity and the required
$C_F\rho^{1-a}$ bound.  Thus every original hypothesis of
BSP10 is supplied, and BSP10 gives
$\texttt{BesselNZ}(\epsilon,\lambda+e_s,u)$.

For the negative shift use \texttt{EvalM}, the dual tuple
$\texttt{nnMinus}(\alpha,-\theta_s,v)$ and its eigenvalue
$-\theta_s-(n-1)$ from BST9.  These are the exact inputs of
\texttt{QuotM}.  Its resulting functions obey the dual half of
BST5, and BSP10 gives
$\texttt{BesselNZ}(\epsilon,\lambda-e_s,u)$.
Both arguments keep every original parameter quantifier and allow
the new existential kernel datum to be the one supplied by BSP10.
Together they say exactly that the original Bessel predicate is
step closed, which is \texttt{StepStatement}.
\end{proof}
% END BST11

\paragraph{What this chapter leaves open.}
The remaining evaluation statements require a nonzero period of the
specified complementary-minor tensor image, with the actual
character and Haar measures.  BST9 does not establish this
nonvanishing; its eigenspace vector may be zero.
The remaining quotient statements must construct $W_f,W'$,
their exact shifted lowest compact type, and period injectivity.
The Casimir computation (BST.35) supplies one required scalar, not
that construction.
Consequently \texttt{StepStatement}, BNV23's unconditional
\texttt{BesselNVGap}, and the subsequent setting remain open at this
stage.  The eventual global-setting join still retains the original
supplied \texttt{Forms0} scope.
